\documentclass[11pt]{article}
\usepackage{amsmath,amssymb,amsthm,mathtools}
\usepackage{enumitem}
\usepackage{geometry}
\geometry{margin=1in}

\title{Step 160: Calkin-Faithful Boundary-Symbol Bridge for \(C_\ell\)}
\author{Six Birds / RH Membrane Working Note}
\date{}

\newtheorem{theorem}{Theorem}
\newtheorem{definition}{Definition}
\newtheorem{proposition}{Proposition}
\newtheorem{nogo}{No-Go Theorem}

\begin{document}
\maketitle

\section{Orientation and typed context}

This track is \emph{adequacy-oriented}.  Smaller residual is better:
\[
\Xi^{\rm BC}\to 0
\quad\text{or}\quad
\Xi^{\rm BC}\preceq \Omega
\]
with a declared budget.

The typed context is the Burnol/Sonine pulled-evaluator carrier.  The comparison
context is the Calkin algebra:
\[
\mathcal Q(H):=\mathcal B(H)/\mathcal K(H),
\qquad q:\mathcal B(H)\to \mathcal Q(H).
\]
There is no untyped comparison between the hard-support model, the Hardy/Hankel
model, and the Burnol/Sonine carrier.

\section{Active residual}

The shifted boundary block is
\[
\mathfrak B_{\ell,a}=J_aP_\infty\tau_\ell(I-P_\infty),
\]
and the Burnol/co-Poisson adequacy residual is
\[
\Xi^{\rm BC}_{\ell,a}
=
\mathfrak B_{\ell,a}^{*}\Pi_{Y_a}\mathfrak B_{\ell,a}.
\]
Burnol's carrier theorem supplies the criterion:
\[
P_a=Y_a^\perp,
\]
so exact adequacy is
\[
\Pi_{Y_a}\mathfrak B_{\ell,a}=0.
\]
After pulling zero-evaluators back to the archimedean/Sonin model, the active
commutator is
\[
C_\ell=(I-P_\infty)M_{m_\ell}P_\infty,
\qquad
m_\ell(s)=e^{-\ell(1/2-s)}.
\]
Let
\[
H_\eta=\overline{\operatorname{span}\{\eta^a_{\rho,k}=J_a^*Y^a_{\rho,k}\}},
\qquad
P_\eta=\operatorname{Proj}_{H_\eta}.
\]
The compact/tail question is
\[
C_\ell P_\eta\stackrel{?}{\in}\mathcal K(H).
\]

\section{Boundary-symbol bridge}

\begin{definition}[Calkin-faithful boundary-symbol bridge]
A Calkin-faithful boundary-symbol bridge for \(C_\ell\) on the pulled-evaluator
context consists of:
\begin{enumerate}[label=(\roman*)]
\item a \(C^*\)-algebra \(\mathcal A_\eta\subseteq\mathcal B(H)\) containing
\(P_\infty,M_{m_\ell},P_\eta\) and the compact ideal \(\mathcal K(H)\cap\mathcal A_\eta\);
\item a boundary-symbol algebra \(\mathcal S_B\);
\item a quotient-compatible \(*\)-homomorphism
\[
\sigma_B:\mathcal A_\eta\to\mathcal S_B
\]
whose kernel on \(\mathcal A_\eta\) is exactly the compact ideal relevant to the
bridge;
\item a boundary leakage operator \(W_\ell\), a lower-faithful partial isometry or
bounded module map \(U_\ell\), and a compact error \(K_\ell\in\mathcal K(H)\) such that
\[
C_\ell=U_\ell W_\ell+K_\ell;
\]
\item a pulled-evaluator inclusion record for \(P_\eta\), so that
\[
\sigma_B(C_\ell P_\eta)=\sigma_B(U_\ell)\sigma_B(W_\ell P_\eta).
\]
\end{enumerate}
The bridge is \emph{Calkin-faithful} if compactness, essential norm, and
Weyl-sequence statements for \(C_\ell P_\eta\) are equivalent to the corresponding
statements for the boundary symbol \(\sigma_B(W_\ell P_\eta)\), up to declared
constants.
\end{definition}

\begin{theorem}[Conditional boundary-symbol equivalence]
Assume a Calkin-faithful boundary-symbol bridge exists and that \(U_\ell\) is
lower-faithful on the boundary range: there are constants \(0<c\le C<\infty\)
such that in the Calkin algebra
\[
c\|q(W_\ell P_\eta)\|\le \|q(U_\ell W_\ell P_\eta)\|
\le C\|q(W_\ell P_\eta)\|.
\]
Then
\[
C_\ell P_\eta\in\mathcal K(H)
\iff
W_\ell P_\eta\in\mathcal K(H),
\]
and more generally
\[
c\,\|q(W_\ell P_\eta)\|
\le
\|q(C_\ell P_\eta)\|
\le
C\,\|q(W_\ell P_\eta)\|.
\]
Thus \(\Xi^{\rm BC}\) is compact/tail-payable exactly when the boundary-Weyl
residual
\[
\Xi^{\rm BW}_{\ell,a}=P_\eta W_\ell^*W_\ell P_\eta
\]
is compact/tail-payable in the bridge context.
\end{theorem}

\begin{proof}
The compact error \(K_\ell\) vanishes under the quotient map \(q\).  Therefore
\[
q(C_\ell P_\eta)=q(U_\ell W_\ell P_\eta).
\]
Lower-faithfulness supplies the two-sided comparison with \(q(W_\ell P_\eta)\).
Vanishing in the Calkin algebra is precisely compactness.  The essential-norm
bounds follow from the same comparison.  Tail-payability then follows from the
standard compact-exhaustion equivalence:
\[
A\in\mathcal K(H)
\iff
\|A(I-Q_N)\|\to0
\]
for every finite-rank exhaustion \(Q_N\uparrow I\).
\end{proof}

\section{Framework-derived no-go: public-shadow non-promotion}

\begin{nogo}[Public-shadow bridge no-go]
Hard-support and Hardy/Hankel computations cannot be promoted to claims about
\(C_\ell P_\eta\) on the Burnol/Sonine carrier unless a Calkin-faithful bridge is
provided.

Scope: this forecloses attacks of the form
\[
\text{Hardy or hard-support model has compactness/noncompactness}
\Rightarrow
C_\ell P_\eta\text{ has compactness/noncompactness}.
\]

Nonclaim: this does not rule out a true boundary-symbol bridge, a direct
co-Poisson factorization, a semilocal prolate theorem, or a direct Burnol/Sonine
Weyl sequence.

Remaining lanes: prove the bridge; construct a pulled-evaluator Weyl sequence;
prove \(C_\ell P_\eta\in\mathcal K\) directly; or prove exact co-Poisson
factorization.
\end{nogo}

\begin{proof}
Compactness is a statement in the quotient \(\mathcal B(H)/\mathcal K(H)\).  A
model calculation on a different carrier supplies no map between quotient
algebras.  Without a quotient-compatible bridge, the model operator and the
Burnol/Sonine operator can have the same finite-window signatures and even the
same formal multiplier while differing by an essential operator.  Therefore a
shadow calculation is non-injective on Calkin data and cannot promote to a
Burnol/Sonine Calkin conclusion.
\end{proof}

\section{Framework-derived no-go: finite-window Calkin blindness}

\begin{nogo}[Finite-window Calkin blindness]
No finite collection of finite-section singular-value audits can certify
\(C_\ell P_\eta\in\mathcal K\) or \(\|C_\ell P_\eta\|_{\rm e}>0\) on the completed
carrier.

Scope: this forecloses finite Jacobi/moment/section evidence as a completed
Calkin certificate.

Nonclaim: an actual completed asymptotic theorem for the finite sections can
still prove compactness or noncompactness.

Remaining lanes: prove a completed limit theorem, construct a Weyl sequence, or
prove a compact decomposition.
\end{nogo}

\begin{proof}
For any finite-rank projection \(Q_N\), the finite matrix
\(Q_N A Q_N\) depends only on a finite corner of \(A\).  One may alter \(A\) on
\(Q_N^\perp H\) by adding either a compact tail or a noncompact essential tail
without changing that finite corner.  Hence finite-window data is not injective on
Calkin classes.  A completed theorem must control the tail uniformly, not just
each finite section.
\end{proof}

\section{Bridge status after Step 160}

The semilocal Connes--Consani--Moscovici framework gives a Hardy--Titchmarsh
and semilocal Sonin ambient geometry.  It supplies the correct room in which a
bridge might live, but the published semilocal results do not, by themselves,
provide the Calkin-faithful boundary symbol for
\[
C_\ell=(I-P_\infty)M_{m_\ell}P_\infty.
\]

Thus the exact missing external theorem is:

\[
\boxed{
\text{Construct }\sigma_B\text{ and prove }
C_\ell=U_\ell W_\ell+K_\ell,\quad K_\ell\in\mathcal K,
\text{ on }H_\eta.
}
\]

The statement of this theorem is framework-derived analytical structural content;
the proof is external semilocal/Sonine operator theory.

\section{Content classification}

\begin{itemize}
\item Predictive structural content: public-shadow promotion and moving-window
support-only verdicts fire in this track; these are applied framework typed
conditions.
\item Analytical structural content: the conditional boundary-symbol equivalence,
the public-shadow bridge no-go, the finite-window Calkin blindness no-go, and
the exact missing boundary-symbol theorem statement.
\item Organizational content: residual tree, bridge atlas, route status, and
nonclaim boundary records.
\item Remaining external content: the proof of the Burnol/Sonine Calkin-faithful
boundary-symbol theorem or the construction of a direct pulled-evaluator Weyl
sequence.
\end{itemize}

\section{Next step}

The next step is Step 161: boundary-symbol theorem extraction or no-go scope
closure.  Target: decide whether existing semilocal Sonin/prolate machinery can
supply the bridge, or whether the Calkin bridge itself becomes the named external
leaf.
\end{document}